\documentclass[conference]{IEEEtran}
\usepackage[utf8]{inputenc}
\usepackage[T1]{fontenc}
\usepackage{cite}
\usepackage{amsmath,amssymb}
\usepackage{graphicx}
\usepackage{booktabs}
\usepackage{url}

\title{Preregistered Paired Evaluation of a Deterministic Verifier Pipeline for LLM‑Generated Network Configuration Changes --- Non‑informative Result under Shared‑Generation Semantics}

\author{
\IEEEauthorblockN{Autonomous AI Researcher}
\IEEEauthorblockA{CNSM 2026 Agentic AI Researcher Track}
}

\begin{document}

\maketitle

\begin{abstract}
We report a fully preregistered paired experiment (cand-2026-01-prereg-v1) that tested whether a deterministic verifier pipeline (generate \ensuremath{\rightarrow} deterministic\_netops\_validator\_v5 \ensuremath{\rightarrow} optional repair) reduces the rate of verifier‑detected unsafe intent\ensuremath{\rightarrow}configuration proposals produced by a hosted LLM on synthetic NetOps tasks. Design: N=200 preregistered unique intent\ensuremath{\rightarrow}configuration tasks in a paired design comparing (a) baseline LLM‑only and (b) guarded pipeline (gpt-5-mini + deterministic\_netops\_validator\_v5). Execution used shared\_initial\_candidate generation semantics (one LLM call per task) producing model\_calls\_used = 200 (preregistered independent per‑arm generation would have used up to 400). Observed (adversarial cluster): baseline SAFE = 200/200; guarded SAFE = 200/200 (paired contingency n11 = 200, n10 = 0, n01 = 0). Estimated paired SAFE‑rate difference = 0.0; exact McNemar two‑sided p = 1.0; paired bootstrap 95\% CI = [0.0, 0.0] (resamples = 10000, seed = 1729). Because identical per‑pair generations were produced and the observed baseline unsafe rate was 0\%, the preregistered causal hypothesis (that the deterministic verifier reduces unsafe proposals) was not testable in this execution. We publish the complete preregistered run and artifacts, provide an artifact‑grounded diagnosis of testability failure, and give concrete, archive‑supported recommendations (independent per‑arm generation budgeting, adversarial injection calibration, and exploratory ROC reserves) for informative repeats. All execution and analysis artifacts cited here are archived in the supplied bundle.
\end{abstract}

\section{Introduction}
Automating human‑intent \ensuremath{\rightarrow} device‑configuration translation with large language models (LLMs) is an active NetOps research thrust because it can reduce human effort and speed maintenance tasks. Yet incorrectly specified or misordered changes can cause outages, policy violations, or security exposures. A commonly proposed mitigation is tool‑augmentation: couple an LLM generator to deterministic validators, sandboxed simulators, or constrained executors in a generate--validate--repair pipeline so that proposed changes are checked before execution.

We preregistered and executed a confirmatory paired experiment (cand-2026-01-prereg-v1) to quantify whether deterministic verifier augmentation causally reduces the frequency of verifier‑detected unsafe proposals from a hosted LLM under both benign and adversarial prompt clusters. The preregistered estimand was the paired SAFE‑rate difference (guarded minus baseline), tested by an exact McNemar two‑sided test with a paired bootstrap CI (10000 resamples, seed = 1729). The design sampled N=200 tasks using a deterministic generator and a paired comparison across two pipelines. Two generation semantics were permitted in the preregistration: independent per‑arm generation (one LLM call per arm per task) or shared\_initial\_candidate (one shared LLM call per task scored by both arms). The executed run used shared\_initial\_candidate to conserve model‑call budget.

Key outcome: under the executed shared semantics the sampled LLM outputs were scored SAFE by the deterministic validator in every confirmatory episode (baseline SAFE = 200/200; guarded SAFE = 200/200). This concordant ceiling outcome (n11 = 200; n10 = n01 = n00 = 0) yields an estimate of zero paired difference and an exact McNemar p = 1.0. Because the observed baseline unsafe rate was 0\% and the same candidate was used in both arms, the preregistered causal hypothesis could not be meaningfully evaluated. The manuscript therefore focuses on three contributions grounded in archived artifacts: (1) a complete, deterministic preregistered run published as reproducible artifacts; (2) a precise, artifact‑grounded diagnosis of why this execution was non‑informative for the confirmatory estimand; and (3) concrete, artifact‑supported design recommendations that would enable informative repeats under the same preregistration framework.

\section{Related Work}
Prior surveys and prototype reports document rapid uptake of LLMs in NetOps workflows (intent translation, configuration synthesis, diagnosis and limited remediation), and they emphasize both promise and safety concerns when models produce executable changes [https://openalex.org/W7161354925; https://openalex.org/W7170714602; https://doi.org/10.36227/techrxiv.177004277.71795055/v1]. Those surveys consistently note that measured correctness gains from LLMs are often task‑dependent and evaluated on curated case sets rather than on replayable workflow‑centred evaluations.

A recurring architectural pattern in the literature is tool‑augmentation: pairing an LLM generator with deterministic tools (validators, simulators, or constrained executors) in generate--validate--refine or divide--verify--refine pipelines. Experimental reports and method proposals show that these pipelines can reduce explicit constraint violations on curated tasks and improve adherence to complex instruction structure when the external verifier/simulator faithfully encodes domain constraints [https://openalex.org/W4412888330; https://openalex.org/W7161354925]. However, the reported improvements depend strongly on the fidelity of the verifier and on the evaluation semantics (how candidate generations are sampled, whether repairs are applied, and whether validation executes in a sandboxed replay), so cross‑study comparisons require careful alignment of these elements.

Work on guardrails, schema‑based tool interfaces, and secure agent orchestration stresses practical trade‑offs: stricter schemas or deterministic checks can lower acceptance of unsafe proposals but may increase false refusals or operator burden if the validator is incomplete or overly conservative [https://doi.org/10.36227/techrxiv.177006109.97957916/v1; https://doi.org/10.36227/techrxiv.177004277.71795055/v1]. Red‑teaming and adversarial evaluation studies from adjacent domains further demonstrate that guardrail designs must be stress‑tested with adversarial templates and that evaluation protocols should report both failure‑to‑detect (false negatives) and false‑refusal (false positive) rates to capture operational trade‑offs; several position and empirical papers advocate ROC‑style diagnostics for verifier behaviour where feasible [https://doi.org/10.2139/ssrn.7132138; https://openalex.org/W7161354925].

Evaluation methodology literature for agentic NetOps emphasizes workflow‑centred, replayable, and sandboxed experiments: tasks should include explicit initial and expected final states, deterministic topology generators, and archived validator traces so that runs are reproducible and auditors can replay failures in a sandboxed environment. Several recent proposals and benchmarks articulate these desiderata and argue that static QA metrics are insufficient when the downstream risk involves stateful multi‑device changes or policy composition across devices [https://doi.org/10.1109/ojcoms.2026.3680457; https://openalex.org/W7161354925].

Our study aligns with this body of work by (a) using a deterministic task generator and sandboxed deterministic validation, (b) preregistering a paired estimand and detailed exclusion/missingness rules, and (c) archiving full validator traces and execution manifests to enable replay. It also addresses an empirical evaluation gap identified in the literature: although tool‑augmentation is widely advocated and shown to help in many reports, there are relatively few reproducible, preregistered adversarial NetOps experiments that quantify the paired causal effect of adding a deterministic verifier under controlled generation semantics. The papers cited above both motivate such paired designs and underscore the sensitivity of conclusions to evaluation semantics (sampling, repair policies, and generation budgeting), which is the methodological locus of our artifact‑grounded diagnosis in this manuscript [https://openalex.org/W4412888330; https://openalex.org/W7161354925; https://doi.org/10.36227/techrxiv.177004277.71795055/v1].

\section{Methodology}
Preregistration and estimand. The experiment was preregistered as cand‑2026‑01‑prereg‑v1. Primary estimand: paired\_success\_rate\_difference\_guarded\_minus\_baseline (guarded - baseline) on binary SAFE indicators obtained from a deterministic sandbox validator. Confirmatory test: exact two‑sided McNemar implemented by paired\_binary\_analysis\_v1; confidence intervals: paired bootstrap percentile (10000 resamples, seed = 1729). The preregistration specified a paired binary design with N=200 tasks, planned paired episodes = 400, and a power plan that targeted detecting an approximate 7 percentage‑point reduction in unsafe proposals with 80\% power assuming a baseline unsafe rate \ensuremath{\approx}10--15\%.

Task generator, scope, and seeds. A deterministic generator (deterministic\_netops\_task\_generator\_v5) produced 200 synthetic intent\ensuremath{\rightarrow}configuration tasks constrained to single‑AS multi‑device topologies and limited configuration fields (admin, mtu, vlan). Generator seeds were deterministic (per‑task generation\_seed 1..200; the RNG seed for the generator family was preregistered) and each task manifest included initial\_state, expected\_final\_state, explicit per‑task workflow policies (required\_sequence), and difficulty metadata. Tasks were clustered by benign vs adversarial prompt transforms using a preregistered obfuscation template set to increase adversarial difficulty; the adversarial cluster was included in confirmatory items per the preregistration.

Pipelines compared. Baseline: accept the LLM candidate (no runtime deterministic enforcement) and score post hoc with deterministic\_netops\_validator\_v5. Guarded: validate the same LLM candidate in a deterministic sandbox using deterministic\_netops\_validator\_v5; the validator emitted structured traces (validation\_before/after, violation lists) and optionally could invoke one automatic repair call per the preregistration. For confirmatory episodes no repairs were applied (repair fields are null in scoring JSONs) so the guarded result reflects pure validation rather than repair‑induced change.

Generation semantics and model budgeting. The preregistered plan allowed two generation semantics: independent per‑arm generation (two model calls per task) or shared\_initial\_candidate (one model call per task whose candidate is scored by both arms). The executed run used shared\_initial\_candidate to conserve hosted model budget. Planned maximum\_model\_calls for independent generation = 400; executed model\_calls\_used = 200. This single shared candidate per task meant that, absent repair, the two arms would always score the same candidate and could not produce discordant paired outcomes arising from arm‑level stochasticity.

Scoring and analysis pipeline. Deterministic\_netops\_validator\_v5 applied full\_intent\_constraint\_validation to each candidate and emitted a deterministic valid flag and detailed violation lists. Confirmatory binary outcomes were derived from validator.valid (1 = SAFE, 0 = UNSAFE). Analysis used paired\_binary\_analysis\_v1 to compute the 2\ensuremath{\times}2 contingency counts (n11,n10,n01,n00), exact McNemar two‑sided p‑value, and paired bootstrap CI (10000 resamples, seed = 1729). Missingness and retry rules from the preregistration were applied: one automatic retry per failed model call; contaminated or nondeterministic tasks are excluded and replaced. Deterministic reconciliation checks and provider audit were recorded as analysis artifacts and used to certify that no technical failures, exclusions, or nondeterminism occurred in the confirmatory run.

\section{Results}
Execution and provider accounting (artifact pointers). The execution manifest records completed\_episode\_count = 400 (200 paired episodes), failed\_episode\_count = 0, and model\_calls\_used = 200 (see execution/execution\_manifest.json and execution/model\_configuration.json for declared model and generation semantics). Deterministic\_reconciliation.json and execution/execution\_manifest.json document that hosted\_model\_call\_event\_count = 200, cache\_miss\_count = 200, and that all deterministic consistency checks passed; no episodes were excluded or flagged (analysis/contamination\_summary.csv reports zero flagged episodes) and analysis/missingness\_summary.csv reports complete\_pairs = 200.

Primary 2\ensuremath{\times}2 paired outcomes and exact inference. The confirmatory contingency table derived by paired\_binary\_analysis\_v1 is: n11 = 200 (SAFE in both arms), n10 = 0 (SAFE guarded, UNSAFE baseline), n01 = 0 (SAFE baseline, UNSAFE guarded), n00 = 0 (UNSAFE both). These counts are archived in analysis/paired\_contingency\_table.csv. From these counts the point estimate for the paired SAFE‑rate difference (guarded - baseline) equals 0.0. The preregistered exact McNemar two‑sided test yields p = 1.0 under the observed contingency. The paired bootstrap percentile 95\% confidence interval computed by the analysis executor (10000 resamples, seed = 1729) is degenerate at [0.0, 0.0]; the analysis log and bootstrap parameters are archived in analysis/analysis\_log.jsonl and analysis/paired\_difference\_figure.svg.

Representative per‑task diagnostics (scoring JSON evidence). Representative scoring artifacts (execution/scoring/task-000001-baseline.json and task-000001-guarded.json) demonstrate that the shared\_initial\_candidate content was identical across arms and that the validator trace fields validation\_before and validation\_after both report valid = true and violation\_count = 0. Repair provider traces are null in all confirmatory scoring JSONs (repair\_provider\_trace\_path = null), confirming that no automatic repair calls altered any candidate during confirmatory episodes. Representative shared\_initial response files (execution/responses/task-000001-shared-initial.txt, task-000002-shared-initial.txt, task-000003-shared-initial.txt) contain the canonical configurations used for scoring; these files match the normalized\_response fields in the scoring JSONs.

Detailed validation metadata. Each scoring JSON embeds validator metadata: scorer\_id = deterministic\_netops\_validator\_v5, scorer\_version = 5.0, scoring\_rule = full\_intent\_constraint\_validation, and scoring\_status = COMPLETED. Validator traces include parsed\_command\_count, required\_command\_count, observed\_touched\_field\_count, and difficulty annotations (e.g., level = medium/hard/very\_hard) for each task; these per‑task difficulty tags are archived in execution/task\_manifest.jsonl and appear in the representative task entries. The validator reported zero violation\_codes and an empty violations array in every confirmatory episode's validation\_after block, consistent with the binary SAFE labels used for the paired analysis.

Why the confirmatory test was uninformative (artifact‑grounded diagnosis). Three concrete archived facts explain the degenerate contingency: (1) the run executed shared\_initial\_candidate semantics (independent\_condition\_generation = false in execution/model\_configuration.json) so both arms scored the same single LLM output per task; (2) the sampled gpt-5-mini outputs were labeled SAFE by deterministic\_netops\_validator\_v5 for every confirmatory task (baseline SAFE = 200/200 as recorded in analysis/condition\_summary.csv); and (3) no repair calls were applied in confirmatory episodes (repair\_provider\_trace fields are null), so the guarded arm could not produce arm‑specific candidate changes. Together these facts produce perfect concordance (n11 = 200) and therefore zero observed discordant pairs, which makes McNemar inference vacuous (p = 1.0) and yields a degenerate bootstrap CI.

Power assumptions vs observed data. The preregistered power calculation assumed a nonzero baseline unsafe rate (\textasciitilde{}10--15\%) and discordant proportions approximated by p10 = 0.07 and p01 = 0.01 (preregistration text), yielding a required sample size \ensuremath{\approx} 174 for 80\% power to detect a \textasciitilde{}7 percentage‑point paired reduction. The executed shared semantics reduced model calls (200 used vs planned 400 for independent per‑arm draws) and removed arm stochasticity; consequently the empirical baseline unsafe rate observed here (0\%) contradicts the power assumptions and collapses inferential sensitivity. This is an execution‑semantic and budgeting outcome rather than evidence about verifier effectiveness.

Exploratory reserves and uncomputed diagnostics. The preregistration reserved up to 50 exploratory tasks to estimate verifier ROC and refusal trade‑offs; that reserve was not consumed before confirmatory execution and therefore no empirical ROC or refusal‑rate estimates were produced in the archived confirmatory artifacts. Analysis/contamination\_summary.csv and analysis/missingness\_summary.csv document that no exploratory ROC artifacts were appended to the confirmatory results. The absence of ROC or refusal data is an archival fact in the supplied bundle and limits secondary claims about verifier true‑positive/false‑negative behavior.

Sensitivity and reproducibility checks. Deterministic\_reconciliation.json shows consistent episode accounting (complete\_pair\_count = 200; n\_11 = 200; provider\_call\_audit.hosted\_model\_call\_event\_count = 200). Analysis artifacts (paired\_contingency\_table.csv, condition\_summary.csv, analysis\_log.jsonl) are self‑consistent and sufficient to reproduce the numeric inference reported here. All raw scoring JSONs, shared\_initial responses, and the model\_configuration/execution\_manifest artifacts are archived in the bundle and permit deterministic replay of the exact shared‑generation execution semantics that produced the degenerate contingency.

\section{Discussion}
We expand the diagnostic interpretation and practical implications of the observed non‑informative (ceiling) outcome, grounding each statement in archived execution and analysis artifacts.

Root cause synthesis. The degenerate contingency (n11=200, n10=n01=n00=0) results from the conjunction of three archived facts: (1) generation semantics: the run executed shared\_initial\_candidate (independent\_condition\_generation = false) and therefore produced exactly one LLM candidate per task that both arms scored; (2) validator outcomes: deterministic\_netops\_validator\_v5 labeled every sampled candidate as valid (baseline SAFE = 200/200); and (3) absence of repairs: repair\_provider\_trace fields are null in confirmatory scoring JSONs, so the guarded arm did not alter any candidate. These archived facts (execution/model\_configuration.json, execution/execution\_manifest.json, execution/scoring/*.json, analysis/paired\_contingency\_table.csv) together make the paired causal contrast mechanically untestable because the arms had no opportunity to disagree.

Implications for confirmatory design. A paired estimand that attributes differences to the presence of a deterministic verifier requires per‑arm variation that could produce discordant outcomes. Under shared generation semantics, discordance can only arise from repair‑driven changes in the guarded arm or nondeterministic validator behavior; neither occurred here. Accordingly, the preregistered power calculation---which assumed a baseline unsafe rate \ensuremath{\approx}10--15\% and discordant proportions that yield detectability with N\ensuremath{\approx}174---no longer applied. The archived model\_call accounting quantifies the specific resource trade‑off: independent per‑arm generation would have required \ensuremath{\approx}400 model calls (preregistered maximum\_model\_calls), whereas the executed shared semantics used 200 calls (execution/execution\_manifest.json). When budgets force a choice, the estimand should drive budgeting: per‑arm independence is necessary to test the causal effect of adding a validator to generation.

Practical recommendations for informative repeats (artifact‑grounded). All recommendations follow directly from the preregistration and execution artifacts:
- Use independent per‑arm generation for confirmatory paired contrasts unless the guarded arm is capable of deterministic repairs that will reliably produce candidate changes. For this study size (N=200), independent generation requires roughly 200 additional model calls (400 total) compared with the executed shared run (200 model\_calls\_used). This arithmetic is visible in the preregistered execution\_contract and the executed manifest.
- Prior to running the confirmatory set, consume a preregistered exploratory reserve (we reserved up to 50 tasks) to calibrate adversarial transforms until the baseline unsafe rate matches the power assumptions. Archive these tuning runs and fix the transform parameters before the confirmatory execution so the confirmatory assumptions remain preregistered and reproducible.
- If repairs are part of the guarded pipeline, instrument, archive, and analyze repair\_provider\_trace artifacts: quantify the fraction of UNSAFE\ensuremath{\rightarrow}SAFE conversions and check whether repairs introduce new violations. In our run repair traces were null; therefore repairs did not serve as a mechanism to induce discordance.
- Implement a deterministic pilot check step that inspects the contingency counts on a small holdout before committing to the full confirmatory test. If the pilot produces a ceiling/floor contingency (no discordant pairs), abort and adapt along preregistered contingency‑preserving paths (e.g., enable independent per‑arm draws or strengthen adversarial transforms) rather than proceeding to a confirmatory test that will be uninformative.

Alternative estimands and complementary analyses. When budgets or policy preclude independent per‑arm draws, consider changing the estimand to one that remains testable under shared semantics: (a) estimate the absolute unsafe‑proposal rate of the LLM baseline (single‑arm proportion) with appropriate CIs; (b) treat the guarded pipeline as an accept/reject gate and estimate refusal and post‑repair SAFE rates (requires repairs and repair traces); or (c) treat the validator as a diagnostic tool and run exploratory ROC estimation (requires ground truth UNSAFE injections and the exploratory reserve). These alternative estimands shift the required resources but preserve inferential validity; the preregistration explicitly listed several secondary\_estimands and exploratory items that map to these alternatives.

Threats to generalization and validator completeness. The deterministic validator used here encodes a formalized full\_intent\_constraint\_validation for the synthetic task family and ran deterministically in the sandboxed environment; however, its completeness relative to production failure modes is limited. The synthetic tasks restrict fields to admin/mtu/vlan and single‑AS topologies; consequently, zero unsafe detections here do not imply general safety in production multi‑vendor or multi‑AS contexts. Future work should validate validators against broader operational rule sets and consider human adjudication for ambiguous cases; such extensions must be preregistered to avoid post‑hoc instrument selection.

Operational implications. From an operator perspective, strict guardrails can reduce risk but may also increase false refusals and operator workload. Our archived run did not produce refusals, so we cannot quantify that trade‑off here; nevertheless, the preregistered exploratory ROC reserve was intended to measure precisely this tension and should be consumed in future repetitions. For deployment decisions, teams should balance the cost of additional model calls (for independent draws or exploratory calibration) against the operational risk of executing unverified changes; our artifacts make the per‑task resource implications explicit and replayable.

Reproducibility and reuse. All artifacts supporting the above diagnosis and recommendations are archived (execution/* and analysis/*). They permit deterministic replay of the executed semantics (shared\_initial\_candidate), verification of model\_call accounting (model\_calls\_used = 200 vs planned 400), and inspection of representative scoring traces that expose why no arm‑level discordance occurred. We encourage future repeats to adopt the pilot‑check and exploratory calibration steps described above and to record the exact decision points in the preregistration so that execution semantics remain auditable and the confirmatory estimand remains testable.

\section{Limitations}
\begin{itemize}
\item Synthetic tasks and scope: tasks were deterministic synthetic topologies produced by deterministic\_netops\_task\_generator\_v5 and limited to single‑AS, multi‑device setups; results may not generalize to production, multi‑AS, or vendor‑specific features.
\item Generation semantics: the executed run used shared\_initial\_candidate (independent\_condition\_generation = false), producing identical per‑pair generations and creating a ceiling effect when shared candidates were valid.
\item Adversarial strength: the preregistered adversarial templates used in this run did not elicit unsafe outputs from gpt‑5‑mini on the sampled tasks; stronger adversarial cases or explicit injected unsafe examples are required to measure verifier effect.
\item Single‑model scope: only the declared model (gpt‑5‑mini) was evaluated; no multi‑model generality is claimed.
\item Verifier completeness: deterministic\_netops\_validator\_v5 ran deterministically in synthetic sandboxed states; external validation of validator completeness against all real‑world NetOps failure modes is not provided.
\item Operational external validity: no production deployment, operator‑in‑the‑loop workload measurement, nor multi‑vendor field trials were performed; operator‑cost trade‑offs remain unquantified at scale.
\end{itemize}

\section{Conclusion}
We executed a fully archived preregistered paired experiment comparing gpt-5-mini baseline generation and a deterministic verifier‑augmented pipeline on 200 synthetic NetOps tasks. Under the executed shared\_initial\_candidate semantics and for the declared model, both arms produced zero verifier‑detected unsafe proposals (paired difference = 0.0; exact McNemar p = 1.0; 95\% CI = [0.0, 0.0]). Because identical per‑pair generations were used and because the observed baseline unsafe rate was 0\%, the preregistered causal hypothesis was not testable in this run. The scientific contribution is therefore methodological and reproducibility‑centred: (1) a complete deterministic preregistered run with archived artifacts that permit exact replay; (2) an artifact‑grounded diagnosis of why this execution was non‑informative for verifier efficacy; and (3) concrete, archive‑supported recommendations (independent per‑arm generation budgeting, adversarial calibration, and exploratory ROC estimation) for informative repeats. The archived execution and analysis artifacts cited throughout (execution/* and analysis/*) permit deterministic replication of the diagnosis and the run outputs and should be consulted prior to attempting repeats or production extrapolation.

\section*{Disclosure Statement}
Disclosure Statement (CNSM Agentic AI Researcher track): Declared LLM: gpt-5-mini (execution recorded in execution/model\_configuration.json). Orchestration/framework: hosted\_netops\_gvr\_v1 adapter and deterministic experiment harness (execution/execution\_manifest.json). Exact initial master prompt: archived at provenance/master\_prompt.txt with immutable SHA‑256 = 1872df1e1805d2d96940456ca016bd665d1d5196add77f5acdf1582bb39b15ba. Topic constraints and autonomy boundary: the human Research Director selected the broad topic family (Generative AI and LLMs for NetOps) prior to the autonomy lock; after the lock the autonomous researcher performed literature review, preregistration (cand-2026-01-prereg-v1), experiment design, code generation, execution, analysis, manuscript drafting, autonomous peer review, and revisions. Preregistration and execution: the preregistration was frozen before execution; key preregistered elements (estimand, paired design N=200, task generator, allowed generation semantics, scoring rules, and stopping rules) are recorded in cand-2026-01-prereg-v1 and archived in the manuscript bundle. Generation semantics and model‑call accounting (executed): independent\_condition\_generation = false (shared\_initial\_candidate semantics) producing a single initial LLM candidate per task shared across arms; model\_calls\_used = 200 while the preregistered maximum\_model\_calls allocated for independent per‑arm generation = 400 (execution/execution\_manifest.json and execution/model\_configuration.json). Validator and scoring: deterministic\_netops\_validator\_v5 performed deterministic sandboxed validation and encoded outcomes as binary labels (1 = SAFE, 0 = UNSAFE) under scoring\_rule = full\_intent\_constraint\_validation; archived scoring JSONs (execution/scoring/*.json) contain validator traces showing validation\_before/after and violation lists. Analysis executor: paired\_binary\_analysis\_v1 computed the paired contingency, exact McNemar two‑sided test, and paired bootstrap CI (resamples = 10000, seed = 1729); analysis artifacts include analysis/paired\_contingency\_table.csv and analysis/condition\_summary.csv. Deviations and restarts: no post‑preregistration design changes were executed; the observed model call count reflects the executed shared\_initial\_candidate semantics. Reproducibility pointers (concise): execution/raw\_results.jsonl, execution/task\_manifest.jsonl, execution/model\_configuration.json, execution/execution\_manifest.json, analysis/paired\_contingency\_table.csv, analysis/condition\_summary.csv, and analysis/paired\_difference\_figure.svg are archived in the supplied bundle and permit deterministic replays. Human editing: no manual human edits to technical content, analyses, or manuscript text occurred after the autonomy lock. Pre‑lock development: infrastructure rehearsals occurred prior to the lock but were excluded from final empirical evidence unless reproduced under the post‑lock execution.

\begin{thebibliography}{99}
\bibitem{ref1} https://arxiv.org/abs/2605.12729
\bibitem{ref2} https://doi.org/10.2139/ssrn.7132138
\bibitem{ref3} https://doi.org/10.36227/techrxiv.177004277.71795055/v1
\bibitem{ref4} https://doi.org/10.36227/techrxiv.177006109.97957916/v1
\bibitem{ref5} https://doi.org/10.1109/ojcoms.2026.3680457
\bibitem{ref6} Xianren Zhang, Xianfeng Tang, Hui Liu, Zongyu Wu, Qi He, Dongwon Lee, Suhang Wang, "Divide-Verify-Refine: Can LLMs Self-align with Complex Instructions?", 2025, 10.18653/v1/2025.findings-acl.709
\bibitem{ref7} Muhammad Bilal, Jon Crowcroft, Ruizhi Wang, Xiaolong Xu, Schahram Dustdar, "Large Language Models for Agentic NetOps and AIOps: Architectures, Evaluation, and Safety", 2026
\end{thebibliography}

\end{document}
